\section{Background and Design Motivation}

\subsection{Why Decentralize LLM Access?}

DIAL is motivated by a simple gap: as LLM services become more diverse, users
still need a way to find and use the right one. A centralized platform can
coordinate discovery, ranking, pricing, and access control, but it also becomes
the place where those decisions are concentrated. This creates several
tensions:

\begin{itemize}
  \item \textbf{Provider concentration}: users often depend on a small number
  of large providers for availability, pricing, and access to strong models.
  \item \textbf{Specialization discovery}: niche model providers need a way to
  be found by clients that need their capabilities.
  \item \textbf{Privacy and locality}: some requests should stay near a user,
  institution, or trusted device.
  \item \textbf{Open participation}: a decentralized service layer should allow
  a new model provider to join without asking a central operator for
  permission.
  \item \textbf{Composable intelligence}: future queries may benefit from
  several models answering, checking, or debating rather than one monolithic
  model responding.
\end{itemize}

The project therefore focuses on one technical layer: service discovery and
routing. The narrower question is whether peers can advertise what they can do,
discover each other, compare profiles, and route a query in a measurable way.

\subsection{Lessons from BitTorrent and Peer-to-Peer Systems}

BitTorrent showed that a decentralized network can coordinate resource exchange
without a central content server \citep{cohen2003bittorrent}. IPFS and libp2p
generalize related ideas: content addressing, peer discovery, transport
abstraction, and protocol negotiation \citep{benet2014ipfs,libp2p}. The useful
analogy is not file sharing itself, but the coordination pattern: a user does
not need to know in advance which machine will provide the resource.

Here, the resource is an LLM service exposed through a node profile. Peers
advertise capabilities such as math, programming, writing, research, planning,
or creative generation. An entry node uses those advertisements to find
candidates and select an executor. The DHT acts as a decentralized capability
index, while direct peer protocols retrieve fuller profiles, request
recommendations, check liveness, and forward queries.

\subsection{Where Blockchain Fits}

Blockchain systems provide shared records that are difficult to modify and a
common order for transactions \citep{nakamoto2008bitcoin}. These properties
can be useful for payments, settlement, provenance, and reputation records.
Blockchain-based AI trust frameworks also use smart contracts, trusted oracles,
and decentralized storage to support auditability and trust
\citep{nassar2020blockchain}. However, putting every capability advertisement
on-chain would impose costly global agreement on a problem that can tolerate
temporary inconsistency.

DIAL therefore separates \textit{discovery} from future
\textit{accounting and trust} layers. Discovery needs low-latency updates: if a
peer leaves, the router can retry or choose another candidate. Payments,
reputation, and provider commitments need stronger integrity and may later
benefit from blockchain or other verifiable records. For the prototype, the
core network is deliberately lighter: libp2p streams for direct communication
and a Kademlia DHT for capability discovery.

\subsection{The PPAI Angle: Matching Supply and Demand}

PPAI helped shape the project vision because it describes a network where
personalized edge LLM agents can be matched to user demand instead of leaving
each user with only one local assistant \citep{wang2026ppai}. The useful idea
for DIAL is the supply-demand framing: users express demand through queries,
while providers express supply through profiles, availability, quality, and
possibly price or reputation.

DIAL narrows that broad idea to a runnable infrastructure layer. Supply becomes
a \texttt{NodeProfile}; demand becomes a small capability vector produced by a
classifier; and the DHT provides decentralized discovery. This is simpler than
PPAI's full matching and
delegation vision, but it gives that vision a concrete base: peers can be
found, compared, selected, and measured before richer policies are added.

\subsection{Design Position}

This positions DIAL as an access-layer prototype for decentralized LLM
services. The implemented focus is peer-to-peer discovery, profile exchange,
and capability-aware routing. Layers such as trust,
reputation, payments, privacy policies, and multi-node execution sit above that
foundation. A DHT does not solve trust or economics, but the network first
needs a way to find services before it can rank, verify, or reward them.
